% texte.tex — f167v, vv. 423–465 (conversion initiale des travaux antérieurs)
\vers[n=423]{}{}{Λακωνίκ᾽ ἄττα, τρεῖς ἔχοντα γομφίους.}{je ne sais quelles laconiennes à trois dents.}
\vers[n=424]{}{}{πρὸ τοῦ μὲν οὖν ἦν ἀλλ᾽ ὑποῖξαι τὴν θύραν}{Autrefois du moins il tenait à nous d'ouvrir furtivement la porte,}
\vers[n=425]{}{}{ποιησαμέναισι δακτύλιον τριωβόλου,}{en nous faisant graver un anneau de trois oboles.}
\vers[n=426]{}{}{νῦν δ᾽ οὗτος αὐτοὺς ᾡκότριψ Εὐριπίδης}{Aujourd'hui ce tâte-poule d'Euripide}
\vers[n=427]{}{}{ἐδίδαξε θριπήδεστ᾽ ἔχειν σφραγίδια}{leur a enseigné l'emploi de cachets vermoulus}
\vers[n=428]{}{}{ἐξαψαμένους. νῦν οὖν ἐμοὶ τούτῳ δοκεῖ}{qu'ils portent suspendus sur eux. Maintenant donc j'opine}
\vers[n=429]{}{}{ὄλεθρόν τιν᾽ ἡμᾶς κυρκανᾶν ἁμωσγέπως,}{que nous combinions de manière ou d'autre un procédé pour le détruire,}
\vers[n=430]{}{}{ἢ φαρμάκοισιν ἢ μιᾷ γέ τῳ τέχνῃ,}{soit par des drogues, soit par un moyen quelconque,}
\vers[n=431]{}{}{ὅπως ἀπολεῖται. ταῦτ᾽ ἐγὼ φανερῶς λέγω,}{afin qu'il meure. Ceci, je le dis à pleine bouche.}
\vers[n=432]{}{}{τὰ δ᾽ ἄλλα μετὰ τῆς γραμματέως συγγράψομαι.}{Le reste, je le coucherai par écrit, avec l'aide de la greffière.}
\vers[n=434, indent=2em, findent=1.2em, metre={lyr.}]{ΧΟ.}{Chœur}{οὔπω ταύτης ἤκουσα}{Jamais je n'entendis}
\vers[n=435, indent=2em, findent=1.2em]{}{}{πολυπλοκωτέρας γυναικὸς}{femme plus féconde en retours}
\vers[n=436, indent=2em, findent=1.2em]{}{}{οὐδὲ δεινότερον λεγούσης.}{ni qui sache mieux parler.}
\vers[n=437, indent=2em, findent=1.2em]{}{}{πάντα γὰρ λέγει δίκαια,}{Tout ce qu'elle dit est juste.}
\vers[suite]{ΧΟ.}{Chœur}{†πάσας δ᾽ ἰδέας ἐξήτασεν†,}{Elle a tout envisagé sous tous égards,}
\vers[n=439, indent=2em, findent=1.2em]{}{}{πάντα δ᾽ ἐβάστασεν φρενὶ πυκνῶς τε}{tout soupesé dans son esprit : et avec discrétion,}
\vers[n=439, indent=2em, findent=1.2em]{}{}{ποικίλους λόγους ἀνηῦρεν}{elle a trouvé un heureux}
\vers[suite]{ΧΟ.}{Chœur}{εὖ διεζητημένους·}{choix d'adroites raisons.}
\vers[n=440, indent=2em, findent=1.2em]{}{}{ὥστ᾽ ἂν εἰ λέγοι παρ᾽ αὐτῆς}{À tel point que si Xénoclès,}
\vers[n=441, indent=2em, findent=1.2em]{}{}{Ξενοκλέης ὁ Καρκίνου, δο-}{fils de Carkinos, parlait tout de suite après elle,}
\vers[suite]{ΧΟ.}{Chœur}{κεῖν ἂν αὐτόν, ὡς ἐγᾦμαι,}{il vous paraîtrait je pense,}
\vers[n=442, indent=2em, findent=1.2em]{}{}{πᾶσιν ὑμῖν}{à tous}
\vers[suite]{ΧΟ.}{Chœur}{ἄντικρυς μηδὲν λέγειν.}{parler sans rien dire.}
\vers[n=443]{ΣΤ.}{Marchande de couronnes}{ὀλίγων ἕνεκα καὐτὴ παρῆλθον ῥημάτων.}{Ce n'est que pour ajouter quelques mots que je m'avance à mon tour.}
\vers[n=444]{}{}{τὰ μὲν γὰρ ἄλλ᾽ αὕτη κατηγόρηκεν εὖ·}{Les autres chefs d'accusation ont été traités à souhait.}
\vers[n=445]{}{}{ἃ δ᾽ ἐγὼ πέπονθα, ταῦτα λέξαι βούλομαι.}{Mais ce que j'ai souffert personnellement, je tiens à le dire.}
\vers[n=446]{}{}{ἐμοὶ γὰρ ἁνὴρ ἀπέθανεν μὲν ἔγκυος}{Mon mari mourut alors que j'étais enceinte,}
\vers[n=447]{}{}{παιδάρια πέντε καταλιπών, ἁγὼ μόλις}{me laissant cinq petits enfants, que je nourrissais à grand'peine}
\vers[n=448]{}{}{στεφανηπλοκοῦσ᾽ ἔβοσκον ἐν ταῖς μυρρίναις.}{en tressant des couronnes sur le marché aux myrtes.}
\vers[n=449]{}{}{τέως μὲν οὖν ἀλλ᾽ ἡμικάκως ἐβοσκόμην·}{En attendant, je gagnais ma vie vaille que vaille.}
\vers[n=450]{}{}{νῦν δ᾽ οὗτος ἐν ταῖσιν τραγῳδίαις ποιῶν}{Mais aujourd'hui cet homme, avec ses tragédies,}
\vers[n=451]{}{}{τοὺς ἄνδρας ἀναπέπεικεν οὐκ εἶναι θεούς·}{a persuadé aux gens qu'il n'y a pas de dieux ;}
\vers[n=452]{}{}{ὥστ᾽ οὐκέτ᾽ ἐμπολῶμεν οὐδ᾽ εἰς ἥμισυ.}{de sorte que notre commerce est réduit environ de moitié.}
\vers[n=453]{}{}{νῦν οὖν ἁπάσαισιν παραινῶ καὶ λέγω}{C'est pourquoi je vous recommande à toutes et vous requiers}
\vers[n=454]{}{}{τοῦτον κολάσαι τὸν ἄνδρα πολλῶν οὕνεκα·}{de le punir, cet homme, pour quantité de raisons.}
\vers[n=455]{}{}{ἄγρια γὰρ ἡμᾶς ὦ γυναῖκες δρᾷ κακά,}{Car il se porte contre nous, ô femmes, à de sauvages violences,}
\vers[n=456]{}{}{ἅτ᾽ ἐν ἀγρίοισι τοῖς λαχάνοις αὐτὸς τραφείς.}{élevé qu'il fut parmi les herbes sauvages.}
\vers[n=457]{}{}{ἀλλ᾽ εἰς ἀγορὰν ἄπειμι· δεῖ γὰρ ἀνδράσιν}{Mais je m'en vas à l'agora.}
\vers[n=458]{}{}{πλέξαι στεφάνους συνθηματιαίους εἴκοσιν.}{Car j'ai vingt couronnes à tresser à prix convenu.}
\vers[n=459, indent=2em, findent=1.2em, metre={lyr.}]{ΧΟ.}{Chœur}{ἕτερον αὖ τι λῆμα τοῦτο}{Cette autre a montré un tour d'esprit}
\vers[n=460, indent=2em, findent=1.2em]{}{}{κομψότερον ἔτ᾽ ἢ τὸ πρότερον}{plus piquant encore}
\vers[suite]{ΧΟ.}{Chœur}{ἀναπέφηνεν.}{que la première.}
\vers[n=461, indent=2em, findent=1.2em]{}{}{οἷα κατεστωμύλατο}{Comme elle a débité,}
\vers[n=462, indent=2em, findent=1.2em]{}{}{οὐκ ἄκαιρα, φρένας ἔχουσα}{non du hors-d'oœvre — elle a du sens}
\vers[n=463, indent=2em, findent=1.2em]{}{}{καὶ πολύπλοκον νόημ᾽, οὐδ᾽}{et un esprit délié — ni de l'abstrus,}
\vers[n=464, indent=2em, findent=1.2em]{}{}{ἀσύνετ᾽ ἀλλὰ πιθανὰ πάντα.}{mais toutes choses emportant conviction !}
\vers[n=465, indent=2em, findent=1.2em]{}{}{δεῖ δὲ ταύτης τῆς ὕβρεως ἡ-}{Oui, il faut que de cet outrage}
\vers[suite]{ΧΟ.}{Chœur}{μῖν τὸν ἄνδρα}{l'auteur}
\vers[suite]{ΧΟ.}{Chœur}{περιφανῶς δοῦναι δίκην.}{subisse un châtiment exemplaire.}
